%=============================================================================!
% 13.1  WHAT A THERMOSTAT MUST DO                                             !
%=============================================================================!
\section{What a thermostat must do}
\label{sec:th-tests}

A real sample keeps its temperature by exchanging energy with what
surrounds it, which a simulation leaves out. A thermostat stands in for
those surroundings. This section sets out what it must achieve, what it
should leave alone, how the chapter tests both, and the single loop into
which every thermostat of the chapter fits.

\subsection*{Two demands}

The first demand is the canonical ensemble of \cref{sec:sm-canonical}. A
run's time averages should equal averages with the weight
$\mathrm{e}^{-\beta H}$, which means three things that can be checked:
the kinetic energy of the $N_{\mathrm f}$ freedoms has the gamma
distribution of \cref{eq:sm-gamma}, with the relative spread
$\sqrt{2/N_{\mathrm f}}$; the positions have the density
$\mathrm{e}^{-\beta U}$ (\cref{sec:sm-canonical}); and the velocities are
Maxwellian and independent of the positions (\cref{sec:sm-maxwell}).
A thermostat that holds the mean temperature but narrows the spread of $K$
samples the wrong ensemble, and every quantity drawn from its
fluctuations, such as the heat capacity of \cref{eq:sm-fluctuation}, is
wrong.

The second demand is softer: the motion should stay close to the motion
of the real system. In a sample the surroundings exchange energy only at
its edges and only slowly, and the atoms inside move by Newton's laws. A
thermostat acts on every atom at every step, so it changes the motion
everywhere, and how much it changes it decides whether a diffusion
coefficient, a rate of hopping or a vibrational frequency measured under
it means anything. The two demands can pull apart: a thermostat that
acts on each atom separately, as those of \cref{sec:th-andersen,sec:th-langevin}
do, disturbs the motion more the more firmly it holds the temperature.

\subsection*{The tests}

Every thermostat in the chapter is tested on two systems. The first is
the liquid argon of Chapter~12, 256 atoms at \SI{135}{\kelvin}, started
from three frames of the run of \cref{sec:sm-ergodic} 20 picoseconds
apart and followed for \SI{100}{\pico\second} with steps of
\SI{10}{\femto\second}, the first \SI{5}{\pico\second} of each set aside
for the distribution of $K$; the three runs give each measured quantity a
mean and a spread. On it we measure the distribution of $K$; the
diffusion coefficient $D$, which measures how fast atoms wander: the
square of the distance an atom moves in a time $t$, averaged over the
atoms and the starting times, grows in proportion to $t$ in a liquid, and
$D$ is a sixth of its slope between 2 and \SI{20}{\pico\second}, with
positions measured from the centre of mass (Chapter~16 treats this
properly, the six included); and the cost, the computing time of a step
against that of velocity Verlet alone.

The second is lithium on the model surface of \cref{sec:en-surface},
standing for lithium in graphite: 1000 atoms, each on its own copy of the
surface with its own thermostat, at \SI{1000}{\kelvin}, started from the
density $\mathrm{e}^{-\beta U}$ and the Maxwell-Boltzmann distribution.
The temperature is high so that hops are frequent: the estimate of
\cref{sec:th-langevin} puts one hop every \SI{1.69}{\nano\second} at
\SI{300}{\kelvin} and \SI{1.26}{\per\pico\second} at \SI{1000}{\kelvin}.
Each atom has only two freedoms, the hardest case for a thermostat: its
kinetic energy should spread as widely as \cref{eq:sm-gamma} with
$N_{\mathrm f} = 2$, a relative spread of 1, and it has no other atoms to
share energy with. Its
exact canonical density of positions is known, so the distribution of
its potential energy can be compared with the exact one, computed on a
grid over one cell of the surface; and its hops can be counted.

\subsection*{The heat and the effective energy}

A thermostat adds or removes kinetic energy at each step. Adding up what
it adds over a run gives the heat it has put in, and the effective energy
\begin{equation}
   E_{\mathrm{eff}} = K + U - \text{heat}
   \label{eq:th-effective}
\end{equation}
counts every change of $K + U$ that the thermostat did not make. Between
the thermostat's actions the motion is Newton's, which keeps $K + U$, so
in exact dynamics $E_{\mathrm{eff}}$ never changes, for every thermostat.
In a run it drifts only through the error of the integrator, and that
drift is the thermostatted counterpart of the energy check of
\cref{sec:cs-trajectory}.

\subsection*{One loop for every thermostat}

Velocity Verlet (\cref{eq:in-kdk}) is a half kick, a drift and a half
kick. Splitting the drift into two halves, as \cref{sec:in-splitting}
split the Liouville operator, gives four pieces, B A A B, with B a half
kick and A a half drift. A thermostat may act before the step, between the
two half drifts, or after the step (\cref{fig:th-step}).
\texttt{mdlab.thermostats.run} is that loop:

\begin{code}
    for step in range(n_steps + 1):
        if step > 0:
            v = act(th.before, v)
            v = v + 0.5 * dt * accel * f
            r = r + 0.5 * dt * v
            v = act(th.middle, v)
            r = r + 0.5 * dt * v
            u, f = model(r)
            v = v + 0.5 * dt * accel * f
            v = act(th.after, v)
\end{code}

\noindent The pass with \texttt{step} 0 makes no step; it and every later
pass end by recording the state, which the listing leaves out.
\texttt{accel} is \texttt{FORCE\_TO\_ACCEL} divided by each mass, and \texttt{act} calls the thermostat's method and adds the change
of kinetic energy it makes to the heat. With no thermostat the three
calls do nothing and the loop is velocity Verlet; the thermostats of the
next sections each fill one or two of the calls. ASE implements five of
them independently. For four, Berendsen's, Langevin's, the Nosé-Hoover
chain and CSVR, from the same start of 32 argon atoms, with the same
random numbers for the two that use them, the two codes agree after 50
steps of \SI{5}{\femto\second} to within \SI{1e-7}{\angstrom} in every
position. ASE's version of Andersen's draws its random numbers
differently, so the two are compared by what they sample: over four runs
of \SI{100}{\pico\second} of the same atoms at \SI{120}{\kelvin}, the
mean kinetic temperature is $(119.40 \pm 0.55)$\,\si{\kelvin} under ASE's
and $(119.95 \pm 0.37)$\,\si{\kelvin} under \texttt{mdlab}'s, each with
0.99 to 1.01 of the canonical spread.

\begin{figure}[htb]
\centering
\begin{tikzpicture}[line cap=round, line join=round,
   box/.style={draw, line width=0.6pt, minimum width=0.9cm,
               minimum height=0.6cm, font=\small},
   th/.style={draw, line width=0.6pt, fill=accent!12, minimum width=1.2cm,
              minimum height=0.6cm, font=\scriptsize, align=center}]
   \node[th] (b0) at (0,0) {before};
   \node[box] (k1) at (1.45,0) {B};
   \node[box] (d1) at (2.6,0) {A};
   \node[th] (mid) at (3.9,0) {middle};
   \node[box] (d2) at (5.2,0) {A};
   \node[box] (k2) at (6.35,0) {B};
   \node[th] (a0) at (7.8,0) {after};
   \foreach \x/\y in {b0/k1, k1/d1, d1/mid, mid/d2, d2/k2, k2/a0}
      \draw[-{Stealth[length=4pt]}, line width=0.5pt] (\x) -- (\y);
   \node[font=\scriptsize, align=center, anchor=north] at (0,-0.5)
      {rescaling\\Berendsen\\CSVR\\chains, $\tfrac12$};
   \node[font=\scriptsize, align=center, anchor=north] at (3.9,-0.5)
      {Langevin};
   \node[font=\scriptsize, align=center, anchor=north] at (7.8,-0.5)
      {Andersen\\chains, $\tfrac12$};
\end{tikzpicture}
\caption{One step of \texttt{mdlab.thermostats.run}: velocity Verlet as a
half kick (B), two half drifts (A) and a half kick, with the places where
each thermostat of the chapter acts; Nosé-Hoover chains act for half a
step before and half after.}
\label{fig:th-step}
\end{figure}

\begin{keyidea}
A thermostat must give the canonical distributions, the gamma
distribution of $K$ with spread $\sqrt{2/N_{\mathrm f}}$ and the density
$\mathrm{e}^{-\beta U}$ of positions, and should disturb the motion as
little as it can. The effective energy, $K + U$ less the heat the
thermostat has added, is kept by the exact dynamics of every thermostat,
and its drift measures the integration error.
\end{keyidea}

\begin{selfcheck}
A thermostat removes \SI{0.5}{\electronvolt} over a run in which $K + U$
falls by \SI{0.48}{\electronvolt}. By how much has the effective energy
changed?
\end{selfcheck}
